% L4-8 (2026-10-30): new subsubsection for §6.5 of numerical_results.tex (label subsubsec:num_L4_8), figure environment included.
% Figure: plots/L4-8/anisotropic_hybridisation.png -> figures/L4_8_anisotropic_hybridisation.png (scripts/export_paper_assets.jl L4-8).
% Data: data/exp_pro/L4-8/{modes,ratios,validation_2d,pair_L40,buffer_absorption,checks}.csv (204/227 checks; failures discussed below).
% Matrix row L4.8 (frozen 2026-10-29, smoke amendment); reduced-model route approved by the PI on 2026-10-30.
{\color{red}
\subsubsection{Anisotropic hybridisation of finite plates}\label{subsubsec:num_L4_8}

\paragraph{Test description.} We test the mode-dilution estimate \eqref{eq:mode_dilution} on free rectangular plates $L_x\times L_y$ with $L_x=20~\si{\meter}$, $L_y\in\{20,40,80\}~\si{\meter}$, $m_\rho=0.01~\si{\meter}$, $C_\rho=0.49~\si{\meter\tothe{5}\per\second\squared}$, $\nu=0$, floating in open water of depth $H\in\{10,50\}~\si{\meter}$, with a uniform homogenised resonator population $n_0m_{\rho,r}=0.01~\si{\meter}$ tuned to each wet mode. The first elastic wet modes of these plates radiate strongly: their radiation damping is $\zeta_\mathrm{rad}=0.08$--$0.22$, while the hybridisation splitting is $\Delta=0.04$--$0.07$, so the two peaks of a forced response merge into one ($\Delta<2\zeta_\mathrm{rad}$, Figure~\ref{fig:num_L4_8}c) and cannot be used to measure $\Delta$. We therefore compute the wet modes by a Ritz reduction of the finite element model. We solve one quarter of the domain (the modes $(1,0)$ and $(0,1)$ are symmetric about both centre lines), with an absorbing frame of width $1.5\lambda$ at a distance $\lambda$ from the plate. The basis is the 25 products $X_i(x)X_j(y)$, $i,j\in\{0,1,3,5,7\}$, of the symmetric dry free--free beam modes. At each frequency, one factorisation of the fluid block of the bare system, with the 25 basis shapes as right-hand sides, gives the projected added-mass matrix $A_B(\omega)$. We interpolate $A_B$ on seven Chebyshev frequencies in $\omega_\mathrm{ref}(1\pm0.12)$ and solve the reduced \eqref{eq:eigWet} by the same fixed point as in \S\ref{subsubsec:num_L2_5}. The homogenised field enters the reduced problem exactly, since $q\propto\eta$ pointwise for a lossless field. Per target mode we obtain $\omega_{mn}$, the splitting $\Delta=(\omega_+-\omega_-)/\omega_{mn}$ with $\omega_r=\omega_{mn}$, the modal mass ratio $\mu_\mathrm{mode}$ and $\zeta_\mathrm{rad}$ from the complex reduced eigenvalue. The fields are $Q_2$, with eight cells per open-water wavelength and eight graded vertical cells ($N=1.1$--$1.8\times10^5$). The same reduction applied to the two-dimensional plates of \S\ref{subsubsec:num_L2_5} reproduces the full eigenvalue computation there to $3.7\times10^{-5}$ in $\omega_1$ and $8.7\times10^{-5}$ in $\Delta$.

\paragraph{Results.} The reduced model represents the finite element solution: the response to the target pressure differs from the monolithic solve by at most $0.63\%$, and a finer mesh (11 cells per wavelength, $N=3.2\times10^5$) changes $\Delta$ by $2\times10^{-4}$. Moving the frame outwards (gap $1.5\lambda$, frame $2\lambda$) changes $A_B$ by $0.2\%$ (real part) and $0.7\%$ (imaginary part), $\Delta$ by $0.65\%$ and $\zeta_\mathrm{rad}$ by $0.56\%$, so the frame also absorbs the obliquely radiated waves. The wet modes mix the dry shapes strongly: the modal assurance criterion of the wet mode against its dry target is $0.40$--$0.82$. On the $40~\si{\meter}$ plate the dry modes $(1,0)$ and $(0,3)$ have close wavenumbers, but the wet problem separates them: the two wet modes are at $2.068$ and $1.915~\si{\radian\per\second}$, with assurance criteria $0.51$ and $7\times10^{-4}$ against $(1,0)$, and $0.07$ and $0.44$ against $(0,3)$. Over all plates the splitting follows $\Delta\simeq\sqrt{\mu_\mathrm{mode}}$ within $2$--$13\%$, and within $15$--$18.5\%$ for the slender $(0,1)$ mode of the $80~\si{\meter}$ plate. The coupling is anisotropic, as predicted: from $\Delta^2$, $\mu_{10}/\mu_{01}=1.72$ and $3.01$ at $L_y/L_x=2$ and $4$ for $H=10~\si{\meter}$, and $1.69$ and $2.73$ for $H=50~\si{\meter}$ (Figure~\ref{fig:num_L4_8}a). The ratio increases with $L_y/L_x$ at both depths. In deep water it grows as $(L_y/L_x)^{0.79}$, against the exponent $1$ of \eqref{eq:mode_dilution}. The modal mass ratios give $\mu_{10}/\mu_{01}=1.90$, $3.51$ ($H=10$) and $1.87$, $3.53$ ($H=50$), closer to $L_y/L_x$. Quantitatively, \eqref{eq:mode_dilution} with $M_f(k_{mn},H)$ overestimates the anisotropy. The deviation is $16\%$ at $L_y/L_x=2$ in deep water, but $27\%$ at $H=10~\si{\meter}$ and $32$--$59\%$ at $L_y/L_x=4$, where the estimate exceeds the bound of $25\%$ set in advance. The local estimate takes the added mass of a single wavenumber $k_{mn}=\beta_1/L$, which is too small for the long, shallow-water $(0,1)$ modes. The absolute splitting is $\Delta^2\approx2\mu_{mn}$, the prefactor already found in two dimensions (\S\ref{subsubsec:num_L2_5}). On the $80~\si{\meter}$ plate, the $(1,0)$ wet mode is no longer identified with the dry shape $X_1(x)$ (assurance criterion $0.14$ and $0.28$), because the $(0,j)$ modes are dense near $k=\beta_1/L_x$. We conclude that a uniform population couples anisotropically to the wet modes of an elongated plate, more strongly to the modes that vary along the short side, as \eqref{eq:mode_dilution} predicts. The estimate gives the trend, but at finite depth and large aspect ratio it overestimates the anisotropy by up to a factor of two.

\begin{figure}[t]
    \centering
    \includegraphics[width=\linewidth]{figures/L4_8_anisotropic_hybridisation.png}
    \caption{Anisotropic hybridisation of free rectangular plates ($L_x=20~\si{\meter}$, quarter domain, Ritz-reduced wet modes). (a) $\mu_{10}/\mu_{01}=\Delta_{10}^2/\Delta_{01}^2$ against $L_y/L_x$ (markers; at $L_y/L_x=1$ the ratio of the diagonal-symmetric and antisymmetric modes), \eqref{eq:mode_dilution} at finite depth (dashed) and the deep-water law $L_y/L_x$ (dotted). (b) Splitting $\Delta$ (circles) and $2\zeta_\mathrm{rad}$ (crosses) against $k_{mn}$, and $\sqrt{\mu_{mn}}$ from \eqref{eq:mode_dilution} (dashed). (c) Response of the reduced model to the target pressure, $L_y=40~\si{\meter}$, $H=10~\si{\meter}$, mode $(1,0)$, without and with the tuned field; dashed: $\omega_\pm$.}
    \label{fig:num_L4_8}
\end{figure}
}
